% ch40-svdq-gemm-w4a4.tex — SVDQuant(三)：W4A4 GEMM（RFC-P3）
% 源码：cache-dit(3295dc6) csrc/kernels/svdq/cute/fp4/gemm_w4a4.{cuh,cu,h}
% 与 ch37 分工：ch37 = LeetCUDA 自研 fp4_gemm.cuh（attention 内外积，non-WS/WS 双路径）；
% 本章 = cache-dit gemm_w4a4（Linear W4A4 + 低秩修正，G2R/TMA hybrid + overlay + LoRA 调度）。
\chapter{SVDQuant（三）：W4A4 GEMM——FP4 MMA 与 BF16 LoRA-up 融合}\label{ch:40}

\section{本章导读}

本章讲 SVDQuant 在线侧的第二个 kernel：消费第~\ref{ch:39}~章 producer 的
三种输出，算出式~\ref{eq:38-svdq} 的完整结果。它是全书最"杂"的一个
GEMM：NVFP4 主积、BF16 低秩修正、两种 TMA 数据源、一条 warp-specialized
流水全在一个 CTA 里。

\textbf{与第~\ref{ch:37}~章的分工}（两章同题于 TMA/WS/SF，读者请对照阅读）：
第~\ref{ch:37}~章讲 LeetCUDA 自研 \texttt{fp4\_gemm.cuh}——attention
内积的量化 GEMM，四路 TMA + non-WS/WS 双路径，重点是 block-scale MMA
的通用流水骨架；本章讲 cache-dit 的 \texttt{gemm\_w4a4}——Linear 的
W4A4 GEMM，\textbf{hybrid 搬运}（A/SFA 走 G2R、B/SFB/D/L1 走 TMA）、
\textbf{stage overlay}（D/L1 与 O 复用 A/B 的 smem 配额）与
\textbf{LoRA 项调度}（LoRA-first/main-first 两种合法时序及其代数补偿），
重点是"一个 accumulator 装两个不同 MMA 的积"。

\textbf{源码区间}：cache-dit commit \texttt{3295dc6} 的
\texttt{csrc/kernels/svdq/cute/fp4/}\texttt{gemm\_w4a4.cuh}（1334 行，
kernel 主体）、\texttt{gemm\_w4a4.cu}（167 行，host 约定与派发）、
\texttt{gemm\_w4a4.h}（71 行，约定注释）；完整代码入口见附录~\ref{app:D}。

\textbf{架构与配置限定}：sm\_120a 及以上（\texttt{\_\_CUDA\_ARCH\_\_ >= 1200}，
非 SM120 分支为 \texttt{trap}）。编译两套 tile 配置：
BN=256/stages=3/BK=128（N、K 对齐允许时的默认；rank tile 32 列，
\textbf{额外要求 rank 是 32 的倍数}）与 BN=128/stages=4/BK=64（窄几何
fallback，rank tile 16 列）。

\textbf{能力边界}（host 校验，\texttt{gemm\_w4a4.cu} L67--137）：
act/wgt 为 packed INT8、ascales/wscales 为 FP8\_E4M3、out 为 BF16；
$M\equiv0\pmod{256}$（v1 激活 pack）、$N\equiv0\pmod{128}$（CUTLASS
128 行 SF atom）、$K\equiv0\pmod{64}$；ascales/wscales 字节数必须精确
等于 \texttt{sf\_buffer\_bytes}；rank $>0$ 时 LoRA 张量为
\texttt{lora\_up[N,R]} BF16、\texttt{lora\_act\_in[M,R]}（layout 决定
FP32 还是 BF16）；\texttt{kNative} 要求 rank $\le128$（定长
\texttt{float[8]} scale 表）。

\textbf{验证方式}：cache-dit 的 \texttt{tests/kernels/test\_svdquant\_cute.py}
与 \texttt{test\_svdquant\_runtime.py}（SM120 门控；本书未复跑）。

\section{问题与动机}

式~\ref{eq:38-svdq} 的求值侧有两个矩阵乘：残差积
$\mathcal Q(\widehat X)\mathcal Q(R)^{\top}$（NVFP4）与低秩修正
$DL_1^{\top}$（BF16）。把它们装进一个 kernel，要同时回答三个工程问题：

\begin{enumerate}[itemsep=0pt,topsep=2pt]
  \item \textbf{两个乘积怎么共享一个 accumulator}？NVFP4 MMA 与 BF16
        MMA 是两种 atom、两种 K 粒度，但输出必须是同一个 FP32 张量
        $s\cdot P+L+\mathrm{bias}$（$s=\alpha\cdot\mathrm{wcscales}$）。
        若两个乘积分属不同 warpgroup，就需要跨 warpgroup 的 FP32 归约；
  \item \textbf{$D$ 与 $L_1$ 怎么进流水而不爆 shared memory}？主循环
        的 stage ring 已占掉 SM120 的 99 KiB 预算大头，再给 LoRA 开
        独立缓冲，stage 数会被迫下降；
  \item \textbf{LoRA 项与主循环谁先谁后}？$D$ 要等 producer 跑完，
        但 $L_1$ 是静态权重；先算 LoRA 可以让它的 TMA 与首个 K tile
        的激活加载重叠，反过来先跑主循环则没有调度风险。
\end{enumerate}

本章的答案是：\textbf{同构 C partition 让两个 MMA 直接共享寄存器
accumulator（问题 1）；D/L1 按字节精确覆盖 A/B 的 stage 配额、输出
tile 最后整片覆盖（问题 2）；两种时序都实现，用一个代数补偿项互换
（问题 3）}。

\section{数学原理（约定层）}

kernel 头注释（\texttt{gemm\_w4a4.cuh} L56--76）把求值约定固定成一套
紧凑的记法，本章全文以它为准。设 $P=\mathcal Q(\widehat X)\mathcal
Q(R)^{\top}$ 为 NVFP4 主积、$L=DL_1^{\top}$ 为低秩修正，输出为
\begin{equation}
  \mathrm{out}=\mathrm{BF16}\bigl(\alpha\,P\,\mathrm{wcscales}
  +\mathrm{bias}+L\bigr),
  \label{eq:40-contract}
\end{equation}
其中 \texttt{alpha} 是标量、\texttt{wcscales} 是逐输出列（$[N]$）的
权重 scale，两者\textbf{只乘主积}。反量化语义按每个 16 元素的 K 组
\texttt{dequant(code[k]) = E2M1(code[k]) * FP8\_E4M3(scale[k/16])}，
与第~\ref{ch:31}~章一致；UE4M3 是\textbf{无符号} e4m3（scale 非负），
CUTLASS 按 128 行 $\times$ 4 factor 的 512 B atom 存储，恰好一个
BK=64 步一盒。物理输入张量沿用 producer 的约定：act 是 $[M,K/2]$ 的
v1 pack（低 nibble 存偶数 K 元素）、wgt 是 $[N,K/2]$ 行主序 pack、
ascales/wscales 都是 FP8\_E4M3。

两条约定决定了后文所有设计。其一，\texttt{alpha} 与 \texttt{wcscales}
\textbf{不能碰 LoRA 项}——低秩分支是高精度修正，乘上量化侧的动态
scale 就破坏了第~\ref{ch:38}~章"低秩分支免费承接大幅值方向"的前提。
其二，两个乘积在同一个 FP32 C fragment 里累加、最后\textbf{统一转
BF16 一次}；这与第~\ref{ch:39}~章 producer 里 $D$ 的单次舍入配对，
构成全链路"每分支一次"的舍入预算。数据流方面的约定同样写在这段
注释里：A/SFA 保持在 producer 输出的 v1 布局里直读，避免运行期转置
或 repack；B 走 shared memory，消除旧实现中每个 warp 重复读权重的
开销；两个 CuTe TiledMMA 的输入类型与 K 长度不同，但 C partition
静态相同——这是下一节"共享 accumulator"的前提。

\textbf{LoRA-first 的代数补偿}。设 $s[n]=\alpha\cdot\mathrm{wcscales}[n]$。
若调度上想让 LoRA 项先算（让它的 TMA/加载与主循环首个 K tile 的激活
预取重叠），则累加器先装 $L=DL_1^{\top}$，而此时残差 scale 还没乘。
kernel 的做法是先把累加器\textbf{除以} $s$（\texttt{apply\_inverse\_scale}），
主积累加后统一乘回并加 bias（\texttt{apply\_residual\_scale}）：
\begin{equation}
  s\cdot\bigl(P + L/s\bigr) + \mathrm{bias} = s\cdot P + L + \mathrm{bias}.
  \label{eq:40-lorafirst}
\end{equation}
补偿的代价分两笔：$\texttt{inv\_scale}[c]=1.0/\texttt{scale}$ 是
\textbf{每列一次除法}（L790--801），随后
$\texttt{tCrC\_rc}(r,c)\;\texttt{*}=\;\texttt{inv\_scale}[c]$ 是
\textbf{每个 accumulator 元素一次乘法}——不是每列一次乘法。换来的
是 LoRA 项与主循环的完全重叠资格。这里没有零保护：若
\texttt{wcscales}$[n]=0$ 或 \texttt{alpha}$=0$，$1/s$ 直接产生
\texttt{inf}/\texttt{NaN} 并污染整个累加器，而 \texttt{main-first}
路径完全不碰这个除法（见 43.8~节常见坑）。bias 不在这一步参与，只由
\texttt{apply\_residual\_scale} 在最后统一加上：\texttt{tCrC} 先由
\texttt{apply\_inverse\_scale} 从 $L$ 变成 $L/s$，主积累加后是
$P+L/s$，收尾乘 $s$ 恰好得到 $s\cdot P+L$。

\section{设计与实现}

\subsection{数据流总览：hybrid G2R + TMA}

实现刻意混用两条输入路径（头注释 L69--72）：
\begin{center}\small
\begin{tabular}{@{}llll@{}}
\toprule
数据 & 路径 & 说明 \\
\midrule
B/SFB（权重+权 scale） & gmem $\to$ TMA $\to$ smem $\to$ RF & 8 个 consumer warp 共享 \\
A/SFA（激活+激活 scale） & v1 pack G2R 直读 & 不过 smem/TMA/mbarrier \\
D/L1（低秩两项） & row-major BF16 $\to$ TMA $\to$ overlay & 复用 A/B 字节配额 \\
\bottomrule
\end{tabular}
\end{center}

\begin{figure}[H]
\centering
\begin{tikzpicture}[font=\scriptsize,
  box/.style={draw, rounded corners=2pt, align=center, inner sep=3pt,
              minimum height=0.95cm},
  arr/.style={->, thick}]
  % 列标题
  \node[font=\tiny, text=tkBlue!80!black] at (2.3,5.2) {gmem 数据源};
  \node[font=\tiny, text=tkBlue!80!black] at (7.9,5.2) {搬运与暂存（producer WG 128T，仅 tid 0 发 TMA）};
  \node[font=\tiny, text=tkBlue!80!black] at (13.2,5.2) {consumer（256T，8 warps，setmaxnreg 224）};
  % 行 1：B/SFB
  \node[box, fill=tkFillBlue, text width=3.3cm] (b) at (2.3,4.2)
    {B/SFB：E2M1 码\\+ ue4m3 SF atom};
  \node[box, fill=tkBlue!18, text width=4.1cm] (sb) at (7.9,4.2)
    {stage ring：$sB+sSFB$\\（S3/S4 轮转）};
  \node[box, fill=tkGreen!25, text width=4.1cm] (ldsm) at (13.2,4.2)
    {LDSM $\to$ $tCrB$/$tCrSFB$};
  \draw[arr] (b) -- node[above,font=\tiny]{TMA} (sb);
  \draw[arr] (sb) -- node[above,font=\tiny]{LDSM} (ldsm);
  % 行 2：A/SFA
  \node[box, fill=tkOrange!22, text width=3.3cm] (a) at (2.3,2.7)
    {A/SFA：v1 256 行\\MMA-fragment pack};
  \node[box, fill=tkOrange!22, text width=4.1cm] (g2r) at (7.9,2.7)
    {G2R 直读（\_\_ldg，\\无 smem/mbarrier）};
  \draw[arr] (a) -- node[above,font=\tiny]{不经 smem} (g2r);
  \node[box, fill=tkFillGreen, text width=4.1cm] (mma) at (13.2,1.5)
    {NVFP4 MMA + BF16 LoRA MMA\\共享 $tCrC$（FP32，单次舍入）};
  \draw[arr] (ldsm) -- (mma);
  \draw[arr] (g2r.east) -- ++(0.55,0) |- (mma.west);
  % 行 3：D/L1
  \node[box, fill=tkPurple!18, text width=3.3cm] (dl) at (2.3,1.1)
    {D $[M,R]$ / L1 $[N,R]$\\BF16 行主序};
  \node[box, fill=tkPurple!25, text width=4.1cm] (ov) at (7.9,1.1)
    {overlay：按字节复用\\A/B stage 配额};
  \draw[arr] (dl) -- node[above,font=\tiny]{TMA} (ov);
  \draw[arr] (ov.east) -- ++(0.55,0) |- (mma.west);
  % 行 4：输出
  \node[box, fill=tkFillBlue, text width=4.1cm] (o) at (13.2,0.1)
    {$O$ tile 覆盖 ring\\STSM $\to$ TMA store};
  \draw[arr] (mma) -- (o);
  % 注记
  \node[anchor=north west, align=left, text width=15.6cm, font=\tiny] at (0,-0.55)
    {三条数据路径互不共享 barrier：主阶段 full 只对 B+SFB 布防（A/SFA 无 TMA 事务），
     LoRA 阶段换成 D+L1 字节数；D/L1 与 $O$ 的复用安全性由 empty 先行到达保证
     （第~\ref{ch:39}~章图~\ref{fig:39-lora-layout} 给出两种 $D$ 形态的对照）。};
\end{tikzpicture}
\caption{gemm\_w4a4 的 hybrid 数据流。B/SFB 与 D/L1 由 producer warpgroup
经 TMA 送入 stage ring（D/L1 按字节覆盖 A/B 配额）；A/SFA 由 consumer
直接从 v1 pack G2R 进寄存器、不与 TMA barrier 交互；两个 MMA 共享
唯一 FP32 accumulator；输出 tile 最后整片覆盖 ring 经 TMA 写出。}
\label{fig:40-hybrid-flow}
\end{figure}

\textbf{为什么是 hybrid 而不是全 TMA}？全 TMA 版本（A/SFB 也走 smem）
在同 shape 的 NCU 对照里 shared 指令多 $22\%$、bank conflict 多
$29\%$、\texttt{math\_pipe\_throttle} 更高（cache-dit 的 NCU 归因，
2026-11-05）；拆掉 A/SFA 的 smem 中转后这些税全部消失。代价是 A 的
\textbf{重复全局 load}：每个 CTA、每个 MMA 线程都按 fragment 位置直读
v1 pack，global 请求是全 TMA 的 4 倍——但这些流量 $87\%$ 命中 L2、
DRAM 吞吐不变，且总量仍只有 native（每 warp 直读全部 fragment）的
$41\%$。hybrid 同时避开"全 TMA 的 smem 税"与"native 的重复 weight
load + spill"，这是对第~\ref{ch:13}~章 TMA 收益边界的一次具体标定：
\textbf{TMA 的价值在共享与解耦，不在本身免费；每个操作数是否需要过
smem，要按它的复用结构与 barrier 成本单独核算}。

\subsection{两个 MMA，一个 accumulator}\label{sec:40-mma}

主积用 Blackwell 的 block-scaled atom；修正项复用 Ampere 时代的
BF16 atom，靠\textbf{同构的 C partition}焊进同一个 accumulator。

先看 tile 几何（源码 L285--301 的选型说明）。一个 CTA 拥有
$128\times$BN 的输出 tile：BM=128 按 32 宽的 N 置换映射成 4 个 M atom
与 2 个 N atom，即 \texttt{AtomLayoutMNK}$=(4,2,1)$，对应 8 个
consumer warp；BK=64 是一个 block-scaled MMA 的 K 长度，BK=128 则在
同一个 smem stage 里放两个独立的 K atom。默认配置是
BN=256/BK=128/stages=3，窄几何回退到 BN=128/BK=64/stages=4（43.4.8
节的派发为此）。MmaAtom 的模板实参为
\texttt{SM120\_16x8x64\_TN\_VS<E2M1, E2M1, F32, UE4M3, 16>}；TN 表示
A、B 都是 K-major，B 的"转置"由 GEMM 语义完成、gmem 中并不转。
LoRA 侧保持同样的 atom layout 与 M/N 置换，只把 K 换成 $kRankTile$，
于是 \texttt{partition\_shape\_C} 逐类型相等，$DL_1^{\top}$ 可以直接
累加进 \texttt{tCrC}，不需要 retile、smem 归约或第二个 accumulator。
两个乘积放在同一批 consumer warp 上是刻意选择：拆成独立的
BF16/NVFP4 warpgroup 会引入跨 warpgroup 的 FP32 归约，还要重复一份
accumulator 存储。下面这段把上述布局固化成编译期常量：

\begin{lstinputlisting}[linerange={332-355},firstnumber=auto,
  title={gemm\_w4a4.cuh L332-355（NVFP4 TiledMMA、LoRA TiledMMA 与一致性 static\_assert）}]{../../../../cache-dit/csrc/kernels/svdq/cute/fp4/gemm_w4a4.cuh}
\end{lstinputlisting}

\texttt{scale\_vec::4X} 的作用是让一条 MMA 指令消费 4 个 FP8 scale
（$4\times16=64$ 个 K），与 BK=64 对齐。两个 \texttt{static\_assert}
（L348--353）把"rank tile 由整数个 BF16 atom 拼成"与"两个 MMA 的 C
partition 一致"在编译期钉死。运行时体现为
\texttt{cute::gemm(lora\_mma, tCrD, tCrL1, tCrC)} 与
\texttt{cute::gemm(tiled\_mma, zip(A,SFA), zip(B,SFB), tCrC)} 写的是
\textbf{同一个} \texttt{tCrC}，无 retile、无 smem 中转、无跨
warpgroup 归约。

\subsection{stage overlay 的内存经济学}\label{sec:40-overlay}

smem 预算的 static\_assert 群（L439--445，其中三条讲 overlay、两条讲
对齐与总预算）是本章最浓缩的工程设计：

\begin{lstinputlisting}[linerange={375-393},firstnumber=auto,
  title={gemm\_w4a4.cuh L375-393（overlay 的字节论证与 D/L1 布局）}]{../../../../cache-dit/csrc/kernels/svdq/cute/fp4/gemm_w4a4.cuh}
\end{lstinputlisting}

\begin{lstinputlisting}[linerange={439-454},firstnumber=auto,
  title={gemm\_w4a4.cuh L439-454（overlay 与预算的 static\_assert + TMA 事务字节）}]{../../../../cache-dit/csrc/kernels/svdq/cute/fp4/gemm_w4a4.cuh}
\end{lstinputlisting}

overlay 的合法性来自位宽的精确互补：BF16 每元素 2 字节，packed E2M1
每元素 0.5 字节，所以一个 BM$\times$(BK/4) 的 BF16 rank tile 与
BM$\times$BK 的 E2M1 stage \textbf{字节数相等}（$BM\cdot BK/4\cdot2
=BM\cdot BK/2$），于是 \texttt{kRankTile = BK/4}（BK=128 时 32 列，
正好一个 SW64 atom 行）。注意 L383--384 的警告：\textbf{字节数匹配
不等于可以复用 A 的 swizzle 布局}——D/L1 有\textbf{自己}的 swizzle
atom（$kRankTile\ge32$ 用 SW64，否则 SW32，L385--389）与独立的 TMA
描述符，只共享存储地址与 stage 步长（L375 与 \texttt{gemm\_w4a4.h} L55
的 "swizzle-free" 是过时注释，按 L385--389 的代码读）；"the empty-stage barrier
proves all consumer warps finished it before L1 overwrites that stage"，
即复用 B 区的安全性由 empty barrier 的先行到达保证。输出 tile
$O$（BF16 SW128 atom）则在 mainloop 与 LoRA MMA 全部结束后整片
覆盖 ring（L443--444 的 static\_assert 要求
$\mathrm{cosize}(O)\le kSmemBytes$）。

代价同样写进了静态断言：BN=256/BK=128 时每 stage 27 KiB（A 8 + B 16 +
SFA 1 + SFB 2），于是 S=2 被输出 overlay（64 KiB $>$ 54 KiB pipeline）
拒绝、S=4 总字节 108 KiB 超 99 KiB 预算——\textbf{stage 数与 overlay 深度
是同一个 99 KiB 配额的两端}，S=3 的 81 KiB 是唯一同时过三条断言的深度；
想加 stage 必须先把 overlay 解耦（把 D/L1/O 挪去独立缓冲），这也是该实现
后续的改造方向。

\begin{figure}[H]
\centering
\begin{tikzpicture}[font=\scriptsize,
  seg/.style={draw, rounded corners=1pt, minimum height=0.75cm,
              inner sep=1.5pt, font=\tiny}]
  \node[anchor=west, font=\small\bfseries] at (0,4.4)
    {同一块 smem 的三次用途};
  % 主阶段
  \node[font=\tiny, anchor=east] at (1.32,3.62) {主阶段};
  \node[seg, fill=tkBlue!22, minimum width=3.9cm] at (3.35,3.62) {A/SFA};
  \node[seg, fill=tkBlue!42, minimum width=5.55cm] at (8.2,3.62) {B/SFB};
  \node[font=\tiny, anchor=west] at (11.15,3.62) {$\times$\,\texttt{kStages}};
  % LoRA 阶段
  \node[font=\tiny, anchor=east] at (1.32,2.32) {LoRA 阶段};
  \node[seg, fill=tkPurple!25, minimum width=3.9cm] at (3.35,2.32)
    {D（BF16 rank tile）};
  \node[seg, fill=tkPurple!42, minimum width=5.55cm] at (8.2,2.32)
    {L1（BF16 rank tile）};
  \node[font=\tiny, anchor=west, text width=4.3cm] at (11.15,2.32)
    {字节数分别等于 A/SFA 与 B/SFB};
  % 输出
  \node[font=\tiny, anchor=east] at (1.32,1.02) {输出};
  \node[seg, fill=tkGreen!28, minimum width=8.9cm] at (6.05,1.02)
    {O tile（BF16 SW128）整片覆盖 ring};
  \node[font=\tiny, anchor=west, text width=4.3cm] at (11.15,1.02)
    {仅当 K 循环与 LoRA MMA 全部完成};
  % 覆盖箭头
  \foreach \x in {3.35,8.2}
    \draw[->, very thick, tkOrange] (\x,3.17) -- (\x,2.75);
  \draw[->, very thick, tkOrange] (5.8,1.87) -- (5.8,1.45);
  % 注记
  \node[anchor=north west, align=left, text width=15.8cm, font=\tiny] at (0,0.35)
    {三条 \texttt{static\_assert} 就是这张图：$\mathrm{kDStageBytes}=\mathrm{kAStageBytes}$、
     $\mathrm{kL1StageBytes}=\mathrm{kBStageBytes}$（BF16 每元素 2 B、E2M1 每元素 0.5 B，
     rank tile 取 BK/4 列，字节数恰好相等）、$\mathrm{cosize}(O)\le\mathrm{kSmemBytes}$。
     D/L1 用自己的 SW64/SW32 atom 与独立 TMA 描述符，只共享地址与 stage 步长；复用安全性
     由 empty barrier 先行到达保证。预算：每 stage 27 KiB，S=4 整环 108 KiB 超 99 KiB 上限；
     S=2 的 ring 只有 54 KiB，装不下 64 KiB 的输出 tile——这正是 S=3（81 KiB）成为唯一
     可行深度的原因。};
\end{tikzpicture}
\caption{stage ring 的三次复用。主阶段存 A/SFA 与 B/SFB；LoRA 阶段用同字节数的 D/L1
覆盖；输出阶段整片覆盖为 $O$ tile。复用靠字节等式与 barrier 顺序，不靠额外缓冲。}
\label{fig:40-overlay}
\end{figure}

SFA 侧还有一个容易忽略的布局细节（L395--413 与 \texttt{filter\_zeros}）：
SF atom 的 K mode 里 16 个码共享 1 字节 scale，layout 含 stride-0
广播；往 atom 追加 stage 维时若按逻辑元素数（$16\times$）推步长会
把下一 stage 写到天上，必须用 \texttt{filter\_zeros} 先剥掉广播模式、
按真实分配字节推进。第~\ref{ch:21}~章 layout 代数的"logical size
与 physical size 分离"在这里是一个会写错内存的实例。

\subsection{warp specialization 与统一流水环}\label{sec:40-ws}

WS 结构（L481--495 注释）：384 线程 = 128 producer + 256 consumer；
只有 producer 的 tid 0 真正发 TMA，其余 127 个保留 warpgroup 身份
参与 \texttt{setmaxnreg}；寄存器预算 producer 32 / consumer 224，
把 384 线程的 CTA 压进 64K 寄存器的 SM（\texttt{warpgroup\_reg\_dealloc<32>}/
\texttt{warpgroup\_reg\_alloc<224>}，L1130/L1155）。

\begin{figure}[H]
\centering
\begin{tikzpicture}[font=\scriptsize]
  \def\cw{2.05}
  \def\topy{2.05}
  \def\boty{0.55}
  % 泳道标签
  \node[anchor=east, align=right, font=\tiny] at (1.32,{\topy+0.38}) {producer tid0\\(TMA)};
  \node[anchor=east, align=right, font=\tiny] at (1.32,{\boty+0.38}) {consumer\\(8 warps)};
  % 阶段标注
  \node[font=\tiny, text=tkPurple!80!black] at (3.2,3.55) {LoRA 阶段（nRt 步）};
  \node[font=\tiny, text=tkBlue!80!black] at (9.5,3.55) {主循环（nKt 步）};
  \draw[thick, tkPurple!70!black] (1.45,3.28) -- (4.95,3.28);
  \draw[thick, tkPurple!70!black] (1.45,3.22) -- (1.45,3.34);
  \draw[thick, tkPurple!70!black] (4.95,3.22) -- (4.95,3.34);
  \draw[thick, tkBlue!70!black] (5.15,3.28) -- (13.35,3.28);
  \draw[thick, tkBlue!70!black] (5.15,3.22) -- (5.15,3.34);
  \draw[thick, tkBlue!70!black] (13.35,3.22) -- (13.35,3.34);
  % 时间刻度
  \foreach \i/\lab in {0/{$g{=}0$}, 2/{$g{=}\mathrm{nRt}$}, 5/{$g{=}\mathrm{gTotal}{-}1$}}
    \node[font=\tiny, text=tkOrange!70!black] at ({1.45+\i*\cw+0.96},3.02) {\lab};
  % producer 泳道
  \foreach \i/\t/\c in {%
      0/{TMA D/L1}/tkGreen!45, 1/{$\cdots$}/tkGreen!18,
      2/{TMA B/SFB}/tkBlue!40, 3/{TMA B/SFB}/tkBlue!40,
      4/{$\cdots$}/tkBlue!18, 5/{TMA B/SFB}/tkBlue!40}{
    \fill[\c] ({1.45+\i*\cw},\topy) rectangle ++(1.92,0.75);
    \node[font=\tiny] at ({1.45+\i*\cw+0.96},{\topy+0.375}) {\t};
  }
  % consumer 泳道
  \foreach \i in {0,1}{
    \fill[tkPurple!30] ({1.45+\i*\cw},\boty) rectangle ++(1.92,0.75);
  }
  \node[font=\tiny] at ({1.45+0.96},{\boty+0.375}) {LoRA MMA};
  \node[font=\tiny] at ({1.45+\cw+0.96},{\boty+0.375}) {$\cdots$};
  \foreach \i in {2,3,4,5}{
    \fill[tkGreen!30] ({1.45+\i*\cw},\boty) rectangle ++(1.92,0.75);
  }
  \node[font=\tiny, align=center] at ({1.45+2*\cw+0.96},{\boty+0.375}) {G2R+LDSM\\+MMA};
  \node[font=\tiny, align=center] at ({1.45+3*\cw+0.96},{\boty+0.375}) {G2R+LDSM\\+MMA};
  \node[font=\tiny] at ({1.45+4*\cw+0.96},{\boty+0.375}) {$\cdots$};
  \node[font=\tiny, align=center] at ({1.45+5*\cw+0.96},{\boty+0.375}) {G2R+LDSM\\+MMA};
  % 泳道间 barrier 说明
  \node[anchor=west, font=\tiny, text=tkOrange!80!black] at (1.45,1.55)
    {每格一对 barrier：producer 按该阶段事务字节 arm full[s]；
     consumer warp leader 完成后 arrive empty[s]，phase $=(g/kStages)\&1$ 在 LoRA 边界不重置};
  % 注记
  \node[anchor=north west, align=left, text width=15.8cm, font=\tiny] at (0,-0.15)
    {LoRA-first 时序（\texttt{kLoRAFirst=true}）：$g$ 从 0 到 $\mathrm{gTotal}{-}1$ 单调推进，
     producer 在 $g<\mathrm{nRt}$ 时装 D/L1、之后装 B/SFB；consumer 先跑 LoRA MMA（此前已预取
     K tile 0 的 A/SFA，故主循环首格跳过重复 load），中间用式~\ref{eq:40-lorafirst} 除 $s$，
     主循环逐 K tile 等待 full[s] 后 G2R A/SFA、LDSM B/SFB、MMA，末段乘回 $s$ 加 bias 并一次性
     转 BF16。rank tile 最多提前 kStages 步进入 ring，消除串行的 LoRA 装载尾巴。};
\end{tikzpicture}
\caption{统一流水环的时序（LoRA-first）。stage 索引与 parity 在 LoRA
边界不重置；full/empty 两类 barrier 分别由 producer 布防字节、consumer
warp leader 到达。}
\label{fig:40-ws-timeline}
\end{figure}

调度上最关键的决定是\textbf{让 LoRA 与主循环共享同一个 step 计数器
g}（L492--495）：nKt 个主 K tile 与 nRt 个 rank tile 排成一个
$[0,\mathrm{nKt+nRt})$ 的全局步序列，stage 索引 $g\bmod kStages$、
parity $(g/kStages)\&1$ 在 LoRA 边界\textbf{不重置}，于是 LoRA 的
TMA 最多可以提前 $kStages$ 步进 ring，与主循环尾部重叠（main-first
模式同理，只是把 rank tile 排在主循环后面）。事务字节数按阶段类型
精确区分（L447--454）：主阶段
\texttt{kTxBytesMain = BN*BK/2 + SFB}（A/SFA 无 TMA 事务），LoRA 阶段
\texttt{kTxBytesLora = D+L1}（\texttt{kNative} 时只有 L1，因为 FP32
$D$ 根本不进 TMA）——注释 L449--450 警告：两类事务字节混在一个
barrier 上会"leave the full barrier waiting forever or release a
stage before its writes complete"。

LoRA 项排在哪里有两种合法选择（\texttt{gemm\_w4a4.cu} L34--40 的
派发）。图~\ref{fig:40-lora-schedule} 把两者的时间线画在一张图里：
\texttt{main-first} 把 rank tile 排在全部 K tile 之后，LoRA 的时间
是一段串行尾巴；\texttt{LoRA-first} 把它排在最前，于是 LoRA 的 TMA
可以与 K tile 0 的 A/SFA 预取重叠。后者需要第~\ref{eq:40-lorafirst}
式的代数补偿，代价只是每列一次乘法。

\begin{figure}[H]
\centering
\begin{tikzpicture}[font=\scriptsize,
  blk/.style={draw, minimum height=0.6cm, inner sep=1.5pt, font=\tiny}]
  % main-first
  \node[font=\tiny, anchor=east] at (1.5,2.25) {\texttt{main-first}};
  \node[blk, fill=tkBlue!30, minimum width=1.5cm] at (2.3,2.25) {K tile 0};
  \node[blk, fill=tkBlue!30, minimum width=1.5cm] at (3.86,2.25) {K tile 1};
  \node[font=\tiny] at (5.1,2.25) {$\cdots$};
  \node[blk, fill=tkBlue!30, minimum width=1.85cm] at (6.6,2.25)
    {K tile $\mathrm{nKt}{-}1$};
  \node[blk, fill=tkPurple!35, minimum width=1.5cm] at (9.0,2.25) {D/L1};
  \node[blk, fill=tkPurple!55, minimum width=1.85cm] at (11.15,2.25) {LoRA MMA};
  \node[font=\tiny, text=tkPurple!70!black, anchor=south] at (10.07,2.62)
    {串行尾巴};
  % LoRA-first
  \node[font=\tiny, anchor=east] at (1.5,1.0) {\texttt{LoRA-first}};
  \node[blk, fill=tkPurple!35, minimum width=1.5cm] at (2.3,1.0) {D/L1};
  \node[blk, fill=tkPurple!55, minimum width=1.85cm] at (4.4,1.0) {LoRA MMA};
  \node[blk, fill=tkBlue!30, minimum width=1.5cm] at (6.6,1.0) {K tile 0};
  \node[blk, fill=tkBlue!30, minimum width=1.5cm] at (8.16,1.0) {K tile 1};
  \node[font=\tiny] at (9.4,1.0) {$\cdots$};
  \node[blk, fill=tkBlue!30, minimum width=1.85cm] at (10.9,1.0)
    {K tile $\mathrm{nKt}{-}1$};
  \node[font=\tiny, text=tkGreen!45!black, anchor=south] at (3.35,1.58)
    {LoRA TMA 与 K tile 0 的 A/SFA 预取重叠};
  % 注记
  \node[anchor=north west, align=left, text width=15.8cm, font=\tiny] at (0,0.35)
    {派发规则（\texttt{gemm\_w4a4.cu} L34--40）：\texttt{kNative} 恒
     \texttt{main-first}（FP32 的 $D$ 不进 TMA stage，只能等主循环结束）；
     \texttt{kRowMajor} 且 $M<16384$ 走 \texttt{LoRA-first}（$M\le4096$ 实测
     $-0.5\%$--$-2.6\%$）；$M\ge16384$ 回退 \texttt{main-first}，因为长序列下
     LoRA 预取会干扰主循环，实测回退约 $+1.4\%$。};
\end{tikzpicture}
\caption{两种 LoRA 调度的时间线。\texttt{main-first} 的 LoRA 是串行尾巴；
\texttt{LoRA-first} 让 LoRA 的装载与首个 K tile 的激活预取重叠，代价是式~\ref{eq:40-lorafirst}
的除法补偿。}
\label{fig:40-lora-schedule}
\end{figure}

\subsection{A/SFA 的 G2R 逆映射}

consumer 侧每 K tile 先 \texttt{wait(full)}，随后逐 BK64 子块执行
"G2R A/SFA $\to$ LDSM B/SFB $\to$ MMA"（L839--858 的主循环）：

\begin{lstinputlisting}[linerange={839-858},firstnumber=auto,
  title={gemm\_w4a4.cuh L839-858（主循环：wait + 逐 k-block 交错）}]{../../../../cache-dit/csrc/kernels/svdq/cute/fp4/gemm_w4a4.cuh}
\end{lstinputlisting}

G2R 本身（\texttt{load\_act\_A\_g2r}，L626--645）是一段手写的 layout
映射：CuTe 的两个 M fragment 对应 $m$ 与 $m+64$，而 v1 pack
每个 256 行块存 8 个 32 行 warp 面板，于是用
\texttt{new\_m/old\_mt/old\_warp/old\_block} 四个量把线程坐标逆映射回
v1 的块地址，SFA 还要再叠一层 \texttt{lane\^{}(2*old\_mt)} 的 XOR。
L1172--1173 的注释把后果说得很重：\textbf{省略这个 swap，形状仍然
对、数学静默错}。这与第~\ref{ch:39}~章 producer 的写出映射是同一
份地图的两端——写端按它排、读端按它解，第~\ref{ch:20}~章的
layout 代数在这里不是抽象审美，是跨 kernel 的接口约定。

\subsection{两条 LoRA 消费路径}\label{sec:40-lora-paths}

\texttt{kRowMajor} 走 \texttt{gemm\_bf16\_lora}（L729--761）：D/L1
经 TMA 进 overlay，LDSM 上寄存器，\texttt{cute::gemm} 直接累加进
\texttt{tCrC}，warp leader 到达 empty。这条路径没有任何逐元素 gather。

\texttt{kNative} 走 \texttt{gemm\_f32d\_lora}（L657--705），差别全在
$D$：FP32 fragment pack 没有 TMA 描述符，consumer 用 CuTe 的
identity-tensor partition（\texttt{partition\_A(cD)}）拿到每个
fragment 值对应的逻辑 $(\mathrm{row},\mathrm{col})$，再逆映射回
native pack 的线性地址逐元素 \texttt{\_\_ldg}，顺手乘上
\texttt{lora\_scales[rank16]} 并转 BF16：

\begin{lstinputlisting}[linerange={685-698},firstnumber=auto,
  title={gemm\_w4a4.cuh L685-698（native FP32 D 的 gather 逆映射）}]{../../../../cache-dit/csrc/kernels/svdq/cute/fp4/gemm_w4a4.cuh}
\end{lstinputlisting}

代价注释（\texttt{gemm\_w4a4.h} L41--45）说得明白：$D$ 是 FP32、
不进 TMA stage，迫使 LoRA 项在 mainloop 之后执行。所以
\texttt{gemm\_w4a4.cu} 的派发（L34--40）是：\texttt{kNative}
$\Rightarrow$ 恒 main-first；\texttt{kRowMajor} 且 $M<16384$
$\Rightarrow$ LoRA-first；$M\ge16384$ $\Rightarrow$ 回退 main-first
（长序列下 LoRA-first 的预取会干扰主循环，实测回退约 $1.4\%$）。

\subsection{epilogue：一次舍入与 TMA store}

\begin{lstinputlisting}[linerange={932-957},firstnumber=auto,
  title={gemm\_w4a4.cuh L932-957（单次 FP32$\to$BF16 + STSM + TMA store）}]{../../../../cache-dit/csrc/kernels/svdq/cute/fp4/gemm_w4a4.cuh}
\end{lstinputlisting}

协议四步（注释 L917--921）：\texttt{\_\_syncthreads()} 证明没有
mainloop stage 还在用这块 smem（overlay 的生命周期到此为止）；
STSM 把 BF16 fragment 写进 swizzled $O$ tile；
\texttt{fence\_view\_async\_shared} 让 generic store 对 TMA 引擎可见、
第二个 \texttt{\_\_syncthreads()} 证明全部 consumer 写完；tid 0 发
TMA store，\texttt{tma\_store\_wait<0>} 排空后才允许 CTA 退出或复用
overlay。

\subsection{host 约定、派发与运行时配对}

host 入口（\texttt{gemm\_w4a4.cu} L58--165）先做完整校验再 launch：
dtype/维度/contiguous/SF 字节数/\texttt{kNative} 的 rank $\le128$/
SM120 能力（L137），然后按几何选 tile 配置——

\begin{lstinputlisting}[linerange={146-163},firstnumber=auto,
  title={gemm\_w4a4.cu L146-163（两套 tile 配置的派发）}]{../../../../cache-dit/csrc/kernels/svdq/cute/fp4/gemm_w4a4.cu}
\end{lstinputlisting}

运行时维度上，Python 侧只有两个公共入口
\texttt{svdq\_gemm\_w4a4} 与 \texttt{svdq\_quantize\_w4a4}，用
\textbf{runtime（native/cute）$\times$ lora\_layout（native/row\_major）}
两个正交维度选择（cache-dit PR \#1259 之后）。配对纪律是硬约束：
\texttt{row\_major} 的 $D$ 必须由本章 producer 体系的
\texttt{quantize\_w4a4\_row\_major} 产出（吃行主序的
\texttt{\_cute\_proj\_down} 与 \texttt{\_cute\_smooth}）；
\texttt{native} 的 $D$ 必须由 native producer 产出（吃 checkpoint 里
fragment-packed 的 \texttt{proj\_down}/\texttt{smooth\_factor}）。
\textbf{混用不报错、结果错 $30$--$54\%$}（cache-dit 实测），这是
"禁止可静默退化的回退链"原则（第~\ref{ch:27}~章）在量化链路上的
现身；cache-dit 用
\texttt{test\_row\_major\_pipeline\_matches\_native\_runtime}
（rank 32/64/128 的端到端 parity）把这对组合钉在回归测试里。

\subsection{forward 时序：从算子到工具}

把三章串起来，一个量化 Linear 层的前向就是两个 kernel：

\begin{center}\small
\begin{tabular}{@{}lll@{}}
\toprule
步骤 & kernel & 产出 \\
\midrule
1 & \texttt{quantize\_w4a4\_row\_major}（第~\ref{ch:39}~章） & qact、ascales（v1 pack）、$D$（BF16）\\
2 & \texttt{gemm\_w4a4(kRowMajor)}（本章） & $Y=\mathrm{BF16}(s\cdot P+L+\mathrm{bias})$ \\
\bottomrule
\end{tabular}
\end{center}

cache-dit 的量化链路（SVDQ PTQ 校准/DQ 无校准、\texttt{runtime\_kernel}
选择、\texttt{cache-dit-convert} 转换 CLI）负责离线产出本链路的权重侧
输入并包装成 \texttt{SVDQW4A4Linear} 模块；本书不展开配置项枚举，
只强调调用形态本身：\textbf{两个 kernel、三次 DRAM 往返以内的在线
侧}（读 $X$、读权重侧、写 $Y$；$D$ 与 pack 驻留显存），就是
第~\ref{ch:38}~章式~\ref{eq:38-svdq} 的全部运行时代价。

\section{图示}

本章的设计图可以按三个问题来读：数据\textbf{从哪来}（图~\ref{fig:40-hybrid-flow}：
A/SFA 直读、B/SFB 与 D/L1 走 TMA）、放在\textbf{哪里}（图~\ref{fig:40-overlay}：
一块 smem 的三种用途）、什么\textbf{时候}算（图~\ref{fig:40-ws-timeline} 的
统一流水环与图~\ref{fig:40-lora-schedule} 的两种 LoRA 排法）。还有一层细节
没进那四幅图：A/SFA 既然不过 smem，consumer 线程就必须知道"我要的 fragment
落在 v1 pack 的哪个字节"——图~\ref{fig:40-g2r-map} 是这个逆映射的三段式。

\begin{figure}[H]
\centering
\begin{tikzpicture}[font=\scriptsize,
  box/.style={draw, rounded corners=2pt, align=center, inner sep=4pt,
              text width=4.3cm, minimum height=1.9cm},
  arr/.style={->, thick}]
  \node[box, fill=tkGreen!18] (cute) at (2.4,1.6)
    {\textbf{CuTe 的 fragment 下标}\\[2pt]
     每个 consumer warp 持 32 行 atom\\
     $\mathrm{mt}\in\{0,1\}$ 是两个 16 行半\\
     $\mathrm{tid}$ 给出 16 行组偏移 $m'$};
  \node[box, fill=tkOrange!18, text width=5.4cm] (inv) at (7.9,1.6)
    {\textbf{逆映射（源码 L626--645）}\\[2pt]
     $\mathrm{old\_block}=m_0/256$，$\mathrm{old\_mt}=(m'/16)\bmod2$\\
     $\mathrm{old\_warp}=(m_0\bmod256)/32$\\
     $\qquad+\,m'/32+2\,\mathrm{mt}$\\
     $\mathrm{lane}'=\mathrm{lane}\oplus(2\,\mathrm{old\_mt})$};
  \node[box, fill=tkBlue!18] (v1) at (13.3,1.6)
    {\textbf{v1 pack 的字节地址}\\[2pt]
     $\mathrm{block\_base}\times1024+\mathrm{old\_mt}\times512$\\
     $\quad+\mathrm{lane}\times16$（字节）\\
     每 lane 一次 \texttt{\_\_ldg} 读 16 B};
  \draw[arr] (cute) -- (inv);
  \draw[arr] (inv) -- (v1);
  \node[font=\tiny, text=tkGreen!55!black, anchor=north] at (2.4,0.35)
    {消费者视角};
  \node[font=\tiny, text=tkOrange!75!black, anchor=north] at (7.9,0.35)
    {换坐标系};
  \node[font=\tiny, text=tkBlue!70!black, anchor=north] at (13.3,0.35)
    {生产者视角};
  \node[anchor=north west, align=left, text width=15.6cm, font=\tiny] at (0,-0.5)
    {三段式的每一段都必须精确对齐（记 $m_0$ 为 CTA 的 128 行块首行、
     $m'=(\mathrm{tid}\bmod128)/32\cdot16$ 为该线程的 16 行组偏移）：
     \texttt{old\_warp} 决定"这个 fragment 属于哪个 32 行面板"，
     \texttt{old\_mt} 决定"面板里哪个 16 行半"，$\mathrm{lane}\oplus(2\,\mathrm{old\_mt})$
     负责修正第二半的 scale lane 交叉。右框的
     $\mathrm{block\_base}=(\lfloor m_0/256\rfloor\cdot K/64+\mathrm{old\_kt})\cdot8+\mathrm{old\_warp}$
     （$\mathrm{old\_kt}$ 为 K64 块号），一个 $\mathrm{block\_base}$ 步进即一块 32 行面板
     （1 KiB）。源码 L1172--1173 的警告是
     这条链上最贵的一句：省略这个 XOR 交换，\textbf{形状仍然对、数学静默错}——错误的数值
     不会崩溃，只会让结果在 PSNR 上静默掉几个 dB。};
\end{tikzpicture}
\caption{A/SFA 的 G2R 逆映射三段式。CuTe 的线程坐标（$\mathrm{tid}$、$\mathrm{mt}$）
先换成 producer 的 warp 面板与 16 行半，再用 lane 的 XOR 修正交叉，最后落到 v1
pack 的字节地址。}
\label{fig:40-g2r-map}
\end{figure}

\section{性能分析}

以下数字全部取自 cache-dit 的测量数据（2026-11-05/06/07），
\textbf{本书未复测}；口径分 CUDA Event（op/kernel 计时）与 nsys
（kernel-only），NCU 仅用于归因、不与 wall-clock 拼接。

\textbf{kernel-only 基线}（RTX PRO 5000，M=4096/N=K=3072，CUDA
Event median，2026-11-05）：rank32 下 cute $120.22\,\mu s$ vs
native $120.24\,\mu s$（持平）；rank128 cute $130.28$ vs native
$132.46$（快 $1.6\%$）。\textbf{rank32 大 shape 仍落后}：
$8192\times6144$ 慢 $4.7\%$、$16\mathrm{K}\times8\mathrm{K}\times
8\mathrm{K}$ 慢 $6.4\%$，归因在 shared 装载量与 TMA/mbarrier
sleeping，是当时的主要瓶颈。

\textbf{PRO 6000 四方对比}（2026-11-07，RTX PRO 6000 Blackwell
SM120，CUDA 13.0.88 / torch 2.10.0+cu130，
\texttt{bench\_svdquant\_kernels.py}，warmup 10 / iters 50，CUDA Event
median，单位 $\mu$s）：表~\ref{tab:40-pro6000} 是 kernel-only GEMM 的
15 格测量，列依次为 native、cute（native producer 的 $D$）、cute
（row-major producer 的 $D$）与 FlashInfer 0.7.3 的 cute-dsl 实现
（后者跑在自己的 $128\times4$ pack 上，cache-dit 各列共享 native
producer 的 pack、row-major 列除外）。

\begin{table}[H]
\centering\small
\caption{RTX PRO 6000 上 SVDQuant W4A4 kernel-only GEMM 的四方对比
（$\mu$s，median；加粗为每行最快）。数据取自 cache-dit 的仓库测量
（PR \#1259），本书未复测。}
\label{tab:40-pro6000}
\begin{tabular}{@{}lrrrrr@{}}
\toprule
$M,N,K$ & rank & native & cute (native $D$) & cute (row-major $D$) & FlashInfer cute-dsl \\
\midrule
$4096,3072,3072$   & 32  & 117.81 & 115.10 & \textbf{112.19} & 148.78 \\
                   & 64  & 121.50 & 118.19 & \textbf{114.00} & 149.58 \\
                   & 128 & 128.40 & 127.90 & \textbf{119.12} & 157.49 \\
$8192,3072,3072$   & 32  & 183.38 & 176.11 & \textbf{174.62} & 205.39 \\
                   & 64  & 188.64 & 183.38 & \textbf{178.10} & 207.14 \\
                   & 128 & 200.03 & 198.72 & \textbf{185.57} & 215.65 \\
$4096,6144,3072$   & 32  & 182.75 & 175.98 & \textbf{173.49} & 201.49 \\
                   & 64  & 189.30 & 184.26 & \textbf{177.23} & 207.90 \\
                   & 128 & 200.50 & 197.26 & \textbf{185.62} & 218.14 \\
$8192,6144,3072$   & 32  & 343.41 & 326.91 & \textbf{310.40} & 361.18 \\
                   & 64  & 337.57 & 321.63 & \textbf{305.90} & 327.41 \\
                   & 128 & 364.38 & 346.16 & \textbf{328.29} & 345.47 \\
$16384,8192,8192$  & 32  & 1590.74 & 1634.34 & 1600.50 & \textbf{1571.23} \\
                   & 64  & 1621.04 & 1654.27 & \textbf{1620.72} & 1658.22 \\
                   & 128 & 1682.77 & 1715.42 & \textbf{1660.46} & 1690.32 \\
\bottomrule
\end{tabular}
\end{table}

读表三条。其一，\texttt{kRowMajor} 的 $D$ 在 15 格中 14 格最快
（相对 native $-0.02\%$--$-9.9\%$），原因是 $D$ 走 TMA 而非 FP32
gather；唯一失手的是 $16384^2$ 的 rank32，被 FlashInfer 以 $1.2\%$
反超，同一格上 cute 的两种 $D$ 都慢于 native。其二，FlashInfer 的
cute-dsl kernel 在 11 格中最慢（$+2.3\%$--$+26.3\%$），说明同一硬件上
"W4A4 算子"的实现差异远大于算法差异——它用的是自己的 $128\times4$
pack，不能与 cache-dit 的 pack 互换。其三，三条 cache-dit 路径数值
一致：对 native 的 \texttt{rel\_l2} 在各种 shape 与 rank 上都是
$3.1$--$3.3\times10^{-3}$（kernel-only 与 e2e 同口径）。

e2e 侧（producer + GEMM）：rank $\le64$ 时
\texttt{cute}+\texttt{row\_major} 8/8 格最优（$-3.7\%$--$-12.6\%$）；
rank $=128$ 时 CuTe producer 退化 $22\%$--$45\%$，e2e 反被 native
反超。dispatch 结论因此是：rank $\le64$ 走
\texttt{cute}+\texttt{row\_major}，rank $=128$ 回退
\texttt{cute}+\texttt{native}。\textbf{该 dispatch 未写成默认}：
PRO 5000 的标定与 6000 有差异，crossover 不能当全局常量。

\textbf{LoRA-first 的收益边界}（2026-11-06，相对 main-first）：
$M\le4096$ 全部受益（rank32--128 约 $-0.5\%$--$-2.6\%$），
$M=16384$ 的 rank32/64 回退约 $+1.4\%$——与 43.4.6~节的派发阈值
$M\ge16384$ 回退 main-first 对应。

\textbf{hybrid 的机制归因}（NCU，2026-11-05）：相对全 TMA，
shared 指令 $-22\%$、wavefront $-18\%$、bank conflict $-29\%$、
静态 FENCE $9\to2$；代价 global 请求 $\times4$（L2 吸收 $87\%$）、
\texttt{long\_scoreboard}$\,2.9$/issue。相对 native，权重走 CTA 内
共享使 global sectors 只有 $22\%$、无 spill（regs 168 vs 252）。

\section{实践经验}

\begin{kaogu}[F4]{gemm\_w4a4 的已证伪路线（cache-dit 的调优结论）}
\begin{itemize}[itemsep=0.5pt,topsep=1pt]
  \item 三方案主循环：A/SFA 前移最大退化仅 $+0.41\%$（未达 $+0.5\%$
        收益门槛）；双缓冲跨 stage 预取 $+1.9\%$--$6.1\%$；BK64 的
        S5--S7 在 2048 上 $+4.6\%$--$7.1\%$、16K 上 $+26\%$--$27\%$。
        并集 $+0.4\%$--$27\%$，全部回退，勿重试；
  \item persistent CTA（grid-stride + epilogue 串行 store）：12 格慢
        $3.0\%$--$6.7\%$；light producer（288T）与 warpgroup 协议不
        兼容（setmaxnreg 死锁），均已回退；
  \item 盲目扫描 stages/tile 已经试过并确认\textbf{无效}：S2 被输出 overlay
        静态拒绝、S4 超 99 KiB；不重测 BN128 小 tile 的旧结论；
  \item direct G2R（全部操作数不过 smem）：对 native pack 仍不可行，
        对 v1 pack 即当前 hybrid——pack 布局决定路线，不是 copy atom
        问题。
\end{itemize}
\end{kaogu}

\section{常见坑}

\begin{itemize}[itemsep=1pt,topsep=2pt]
  \item \textbf{full barrier 字节数混配}：主阶段与 LoRA 阶段的
        \texttt{expect\_tx} 字节数不同，混用会死等或提前放行；
  \item \textbf{overlay 复用布局而非仅字节}：BF16 D/L1 不能套 A/B 的
        SW32 swizzle 布局，只能共享地址与 stage 步长（L383--384）；
  \item \textbf{G2R 逆映射的 lane XOR}：SFA 的第二 m-tile 需要
        \texttt{lane\^{}2}，省略则形状对、数值错；
  \item \textbf{SF stage 步长的 filter\_zeros}：stride-0 广播模式下
        按逻辑 size 推进 stage 会写飞，先剥广播再取真实字节；
  \item \textbf{e2e PSNR 判别不了 kernel 等价}：SVDQ 路径同配置两次
        出图仅 $34$--$36\,\mathrm{dB}$，与跨实现差异同量级；kernel
        等价性只能看 parity 用例（\texttt{test\_svdquant\_cute.py}）；
  \item \textbf{kernel-only 高估收益}：quantize 占 e2e 的 $14\%$--$31\%$
        且随 rank 上升，评估收益必须回到完整 Linear 层；
  \item \textbf{rank128 不要用 row\_major producer}：PRO 6000 上退化
        $22\%$--$45\%$，与 layout 的配对纪律一起看；
  \item \textbf{LoRA-first 的 $1/s$ 没有零保护}：
        \texttt{apply\_inverse\_scale}（L790--795）直接算
        $1.0/\texttt{scale}$，\texttt{wcscales} 某列为 0 或
        \texttt{alpha}$=0$ 时产生 \texttt{inf}/\texttt{NaN} 并污染整个
        FP32 累加器；\texttt{main-first} 路径不碰这个除法，校准阶段若
        出现全零列可以切过去绕过。
\end{itemize}

\section{面试要点速查}

\begin{itemize}[itemsep=1pt,topsep=2pt]
  \item Q：两个不同 MMA（NVFP4/BF16）怎么共享 accumulator？\\
        A：AtomLayout 与 M/N permutation 同构使 \texttt{partition\_shape\_C}
        逐类型相等（static\_assert），\texttt{cute::gemm} 直接写同一
        \texttt{tCrC}（\ref{sec:40-mma}~节）。
  \item Q：D/L1 复用 A/B 的 stage 配额为什么合法？\\
        A：BF16 2B/elem vs E2M1 0.5B/elem，BM$\times$(BK/4) 的 BF16
        tile 与 BM$\times$BK 的 E2M1 stage 字节数相等；empty barrier
        保证消费者先撤（\ref{sec:40-overlay}~节）。
  \item Q：LoRA-first 的代数补偿是什么？\\
        A：先算 $L$ 时把累加器除以 $s=\alpha\cdot\mathrm{wcscales}$，
        主积后统一乘回：$s(P+L/s)+\mathrm{bias}=sP+L+\mathrm{bias}$
        （式~\ref{eq:40-lorafirst}）。
  \item Q：为什么 A/SFA 不走 TMA？\\
        A：v1 pack 按 fragment 交织，G2R 直读免 smem 中转与 barrier，
        shared 指令 $-22\%$；代价是重复 load，由 L2 $87\%$ 命中吸收。
  \item Q：kNative 与 kRowMajor 的 $D$ 路径差在哪？\\
        A：FP32 fragment + gather + 在线 scale，且强制 main-first；
        vs BF16 行主序 + TMA + 可 LoRA-first（图~\ref{fig:39-lora-layout}）。
  \item Q：这个 kernel 的 stage 数为什么卡在 3？\\
        A：99 KiB 预算内，S=4 总字节 108 KiB 超限，S=2 的 54 KiB ring
        装不下 64 KiB 的输出 tile；加 stage 须先解耦 overlay。
  \item Q：producer 与 layout 配对的坑？\\
        A：\texttt{row\_major} $D$ 只配 \texttt{\_cute\_proj\_down}/
        \texttt{\_cute\_smooth}，混用静默错 $30$--$54\%$，用端到端
        parity 测试钉死。
\end{itemize}

\section{最小测试与动手实验}

\begin{lstlisting}[style=console]
# 读者侧验证（cache-dit 仓库，SM120 卡；本书未复跑）：
$ python -m pytest tests/kernels/test_svdquant_cute.py \
                   tests/kernels/test_svdquant_runtime.py -q
# 上面两个文件覆盖 producer 数值、rank 32/64/128 的 GEMM parity、
#   runtime x lora_layout 四组合与 INT4 隔离回归；cache-dit 记录的
#   94 个用例是 5 个测试文件的合计（另含 ops 派发、source manifest、
#   quantize config）

# CPU 侧观察（layout 代数复算，PyCuTe 纯 host）：
# 用 PyCuTe 复算 SmemLayoutA(SW32 atom, (BM,BK,stages)) 与 SmemLayoutD
# (SW64 atom, (BM,BK/4,stages)) 的 cosize，验证
#   cosize(D)*2B == cosize(A)/2*1B == kAStageBytes   （overlay 字节恒等式）
# 以及 sf_gmem_layout 的 512B atom 内偏移公式（L89-98）。
\end{lstlisting}

\section*{延伸阅读}

\begin{itemize}[itemsep=1pt,topsep=2pt]
  \item cache-dit \texttt{bench/kernels/bench\_svdquant\_kernels.py}：
        本章全部性能数字的复测入口（SM120；本书未执行）；
  \item CUTLASS \texttt{Sm1xxBlockScaledConfig}：512 B SF atom 的
        布局定义与本章 \texttt{SfSmemAtom} 的对照；
  \item nunchaku 的 W4A4 kernel：同一数学约定的上游实现，G2R/TMA
        分工与本章相反（权重 fragment 直读），对照阅读可体会
        "谁过 smem"的权衡空间；
  \item 本部分另两章：第~\ref{ch:38}~章（数学主线）、第~\ref{ch:39}~章
        （本章 GEMM 的上游 producer）。
\end{itemize}
